% TOP_PROBLEM: {"id": 4800029, "title": "Multiple ergodic averages — Problem 29", "category_id": 11, "category": {"id": 11, "name": "dynamical_systems", "display_name": "Dynamical Systems", "description": "Problems about long-term behavior of deterministic systems, Hamiltonian dynamics, and periodic orbits.", "slug": "dynamical-systems", "order_index": 11, "created_at": "2026-07-31T15:26:25.671Z"}, "status": "partially_solved", "problem_number": "AMR-047-0029", "set_id": 13}
% FIRST_POSTED: 2026-10-01
% ENTRY_KIND: statement_only
% UnsolvedMath ID 4800029; source catalogue code AMR-047-0029
% !TeX program = lualatex
% Definition and sources only; this file is not an LLM solution attempt.
\documentclass[11pt,a4paper]{article}
\usepackage[margin=25mm]{geometry}
\usepackage{fontspec}
\setmainfont{lmroman10-regular.otf}[BoldFont=lmroman10-bold.otf,
 ItalicFont=lmroman10-italic.otf,BoldItalicFont=lmroman10-bolditalic.otf]
\usepackage{amsmath,amssymb,mathtools}
\usepackage{unicode-math}
\setmathfont{latinmodern-math.otf}
\AtBeginDocument{\let\setminus\smallsetminus}
\usepackage{xurl,hyperref}
\hypersetup{colorlinks=true,urlcolor=blue}
\setcounter{secnumdepth}{0}
\setlength{\parindent}{0pt}
\setlength{\parskip}{5pt}
\setlength{\emergencystretch}{3em}
\begin{document}
\section*{Multiple ergodic averages — Problem 29}\label{4800029}
{\small UnsolvedMath ID 4800029 · AMR-047-0029 · Dynamical Systems · AMR Open Problem Lists}
\subsection{Definitions and mathematical statement}
Fix a real number $c>1$ which is not an integer,
and put $a(n)=\lfloor n^c\rfloor$ for $n\ge1$.
Let $T_1,\ldots,T_\ell$ be commuting invertible
measure-preserving maps on a probability space
$(X,\mathcal B,\mu)$, with arbitrary $\ell\ge1$.
For bounded measurable $f_1,\ldots,f_\ell$, define
\[
 A_N=\frac1N\sum_{n=1}^N\prod_{i=1}^\ell f_i\circ T_i^{a(n)},
 \qquad
 B_N=\frac1N\sum_{n=1}^N\prod_{i=1}^\ell f_i\circ T_i^n.
\]
Prove that $A_N$ converges in $L^2(\mu)$ and
that its limit equals the $L^2$ limit of $B_N$.
Equivalently, in addition to existence of the
ordinary mean limit, prove $\|A_N-B_N\|_2\to0$.

The associated recurrence assertion is that every
measurable $A$ of positive measure admits some
$n\ge1$ for which
\[
       \mu\bigl(A\cap T_1^{a(n)}A\cap\cdots
                             \cap T_\ell^{a(n)}A\bigr)>0.
\]
All systems and all finite families are included,
without ergodicity, mixing, or independence
assumptions. The limit need not be a product
of integrals; it is the ordinary commuting
average's own limit. The exponent $c$ is fixed
before the limit in $N$ is taken, with no
uniformity as $c$ varies required. Integer powers
$c$ are excluded from this particular assertion.
\subsection{Short English statement}
Do noninteger power times give the same mean limits and the same universal recurrence as ordinary times for commuting transformations?
\subsection{Sources}
\begin{itemize}
\item[S1] ULAM.AI and the UnsolvedMath Contributors, supplied problems.json, record 4800029 (AMR-047-0029), AMR Open Problem Lists. Catalogue identity and source transcription; CC BY 4.0.
\url{https://huggingface.co/datasets/ulamai/UnsolvedMath}.
\item[S2] Frantzikinakis - Some open problems on multiple ergodic averages (2016), Problem 29. Original source indexed by AMR-047-0029.
\url{https://arxiv.org/abs/1103.3808}.
\end{itemize}
\end{document}
